\label{section:geometry} The advection equation for the volume fraction field
of a phase $\alpha$ for an incompressible flow is written in the following
form:
%
\begin{equation}
  \partial_t{\alpha} + \nabla \cdot {(\U \alpha)} = 0. 
  \label{equation:alpha-advection}
\end{equation}
%
Applying the discretization practice of the unstructured FVM, with the first
order Euler explicit time discretization on Equation
\ref{equation:alpha-advection}, results with the following algebraic equation: 
%
\begin{equation}
  \alpha_{c}^n = \alpha_{c}^o - \frac{\Delta t}{V_c} 
  \sum_f  \alpha_{f} \vec{U}_f \cdot \vec{S}_f, 
  \label{equation:alpha-discrete}
\end{equation}
%
where $f$ denotes the face of a computational cell (cp. Fig.\ 
\ref{fig:polyhedralCell-FVM}) and, correspondingly, $\vec{U}_f$ and 
$\vec{S}_f$ are the face-centered velocity and face area normal 
vectors shown in Figure \ref{fig:polyhedralCell-FVM}. 

\begin{figure}[h]
	\centering
    \def\svgwidth{0.4\columnwidth}
       {\footnotesize
        \input{figures/polyhedralCell-FVM.pdf-tex}
       }
	\caption{Polyhedral cell}
	\label{fig:polyhedralCell-FVM}
\end{figure}

The unstructured FVM utilises a computational domain decomposed into
non-overlapping convex polyhedra (Figure \ref{fig:polyhedralCell-FVM}) which are the 
building elements of the unstructured mesh. However, the implementation of the mesh 
is done in an efficient way, when it comes to addressing the information needed by 
numerical algorithms. More precisely, the mesh points are defined only once as a 
set of position vectors, and are indirectly addressed by the unstructured FVM 
algorithms. The indirect addressing leads to the definitions of faces and cells 
as sets of indirect indices: the face is thus an ordered set of indices which 
relate (map) to the list of mesh points, and the mesh cell is a set of indices 
which relate to the list of mesh faces. 

The ordering of the face indices determines the orientation of the face area
normal vector $\vec{S}_f$, shown in Figure \ref{fig:polyhedralCell-FVM}. For
the polyhedron in Figure \ref{fig:polyhedralCell-FVM}, the cell index in the
list of mesh cells is lower than the cell index of its neighbour accross the
face $f$. This makes the polyhedron the \emph{owner} of the face $f$, and makes
the ordering of the indices of the face $f$ such that the face area normal vector
$\vec{S}_f$ points outwards from the owner polyhedron, and into the neighbor polyhedron.
Such ordering allows for efficient (unique) computation of the volumetric
fluxes used by unstructured FVM. 

An algebraic approach to solving Equation \ref{equation:alpha-discrete}
utilizes bounded higher-order interface capturing (compressive) advection 
schemes on unstructured meshes to determine the face-centered volume fraction 
value $\alpha_f$. This practice often results in an artificial smearing of 
the volume fraction field as well as artificial deformation of the interface. 

The geometrical VoF algorithm solves Equation \ref{equation:alpha-discrete} by
computing the fluxed phase volume: 
%
\begin{equation}
    V_{f,\alpha} = \alpha_{f} \| \U_f \cdot \vec{S} \| \Delta t = \alpha_f \| F_f \| \Delta t
    \label{equation:vol-phase-flux}
\end{equation}
%
using geometrical operations, where $F_f$ is the volumetric flux across the face
$f$ with an outward pointing surface area normal vector $\vec{S}$ 
($\vec{S}_f$ for a computational cell formally owning the face $f$ -- cf.\ Figure 
\ref{fig:polyhedralCell-FVM}), which satisfies the discrete divergence free condition: 
%
\begin{equation}
    \sum_f F_f = 0.
    \label{eq:discretedivergence}
\end{equation}
%
Without the indirect addressing of the unstructured mesh, \eqref{eq:discretedivergence} 
would be computed on a per-cell basis, repeating the flux calculation twice per cell 
face. To avoid this, the flux is computed using the oriented face area normal vector 
$\vec{S}_f$ and the owner-neighbor addressing of the unstructured mesh: 
%
\begin{equation}
   \sum_f \alpha_f \vec{U_f}\cdot\vec{S} = \sum_{\text{owner}} \alpha_f \vec{U_f}\cdot\vec{S}_f 
     - \sum_{\text{neighbor}} \alpha_f \vec{U_f}\cdot\vec{S}_f, 
   \label{equation:own-nei}
\end{equation}
%
thus computing the same flux once per cell and adding the flux contribution to
the owner cell (outward directed $\vec{S}_f$) and deducting the flux
contribution from the neighbor cell (inward directed $\vec{S}_f$). 

Because of the geometrical multidimensionality of the directionally
unsplit advection algorithm, the artificial deformation of the interface which
is transported in the \emph{skew direction}\footnote{Skew direction in the mesh
is the direction of the point or edge neighbor cells.} is significantly
reduced, compared to fluxed-based methods based on algebraic differencing schemes. 
Such geometrical way of computing the fluxed phase volume also ensures the 
interface geometry to remain located within a single layer of cells throughout 
the simulation which is consistent with the discrete counter-part of the underlying 
sharp interface model. 

The geometrical VoF method relies on geometrical information of the interface. 
In our work, we have used a standard approach for the geometrical VoF
algorithm: a planar geometry of the interface is placed within an
\emph{interface cell} and the geometrical evolution of the volume fraction
field is performed by two subsequent steps, namely the \emph{interface
reconstruction} and the \emph{volume fraction advection} step. This algorithm
is commonly known as Piecewise Linear Interface Calculation (PLIC) algorithm
\citep{RiderReconstructing1998}.

In the remainder we provide a detailed explanation of the supporting
geometrical operations and both geometrical steps of our geometrical VoF
method. 

\subsection{Basic geometrical operations}

Two geometrical operations are the basis for both the reconstruction and the
advection step: the intersection between a plane and a polyhedron, and the
intersection between two polyhedrons. These operations belong to the field of
\emph{Computational Geometry} \citep{SchneiderGTC, DeBerg2008CG} and are
usually sources of difficulties when it comes to calculations that support
geometrical advection of the volume fraction in three dimensions. The reason
for this lies in the fact that the polyhedrons involved in the intersection
operations become non-convex due to the local complexity of the discrete
velocity field. When the full Lagragian discrete flow map is used to compute
the flux polyhedrons, complex flow fields may introduce non planarity of the
swept faces. As reported by \citet{AhnShashkov2009}, in case of a full
Lagrangian backtracking of the whole mesh, and subsequent re-mapping onto the
Eulerian static mesh, the faces of the swept cells may be non-planar, and the
resulting polyhedrons non-convex. If this occurs, a simplification of the swept
mesh is performed that enforces geometrical operations on convex polyhedrons
for the sake of simplicity and robustness of the overall algorithm.

Intersection between a plane and a convex polyhedron is calculated using the
\emph{polyhedron clipping and capping algorithm}, using a slightly modified
version of the algorithm described by \citet{AhnShashkovGeometricAlg2008}. As 
the method's name suggests, each face of the polyhedron is visited and 
\emph{clipped} by a plane, marking the new vertices that will at the end be 
added as the polyhedron \emph{cap} face. Once all the faces are clipped, the 
capping algorithm will append the final cap face making sure there are no 
repeated points. The polyhedron faces are defined to be oriented counter 
clockwise, when viewed from outside of the polyhedron. Special cases may 
appear during the geometrical operation, which may result in an intersection 
being represented as: void, point, edge, face or the complete polyhedron. 
All these cases are internally and automatically dealt with within the 
clipping part of the algorithm implementation.

\begin{figure}
    \centering
    \def\svgwidth{0.5\columnwidth}
       {\footnotesize
        \input{figures/clip-cap-black-white-scaled.pdf-tex}
       }
    \caption{Cutting (clippig and capping) of a polyhedral cell}
    \label{fig:cap}
\end{figure}

The final cap face of the polyhedron is shown in Figure \ref{fig:cap}. The
order of the cap vertices $p_{cp}$ could be computed using the connectivity of
the clipped polyhedron: cap vertices belong to edges that have only a single
face attached to them after the clipping part of the algorithm is finished. In
Figure \ref{fig:cap}, $cp$ marks the cap point and $p_{c}$ is the center
point of the cap. However, we have sorted the points based on the angle the
vertex makes with respect to the cap center axis, given by the first cap point
candidate. Since we are dealing with polyhedra that have small number of faces
(finite volume cells usually have a small number of faces), an angle based
insertion sort is less complicated to implement and is efficient when compared
to constructing the face connectivity information to locate cap edges of the
polyhedron using a graph representation of the polyhedron. In order to find the
cap edges using edge-face connectivity, an edge based graph representation of
the polyhedron would be necessary. Once the clipping and capping of the
polyhedron is performed, the intersection between the polyhedron and a plane is
complete. 

The intersection between two convex polyhedrons can be built by applying the
polyhedron clipping and capping algorithm for each face of a polyhedron. The
same algorithm is used for calculating the intersection between a plane and a
cell, additionally ensuring the face normal is directed outwards from the
polyhedron created from the mesh cell in OpenFOAM. 

\subsection{Interface reconstruction}

We have chosen a first order accurate reconstruction algorithm similar to the
one of \citet{ParkerYoungsGradient} in which the gradient of the volume
fraction field determines the orientation of the interface plane with a
difference in the accurate and parallelized method of gradient calculation on
arbitrary unstructured meshes. This algorithm is first order convergent
regarding the geometrical reconstruction error, i.e. it cannot reconstruct
planar interfaces exactly. The absolute accuracy in determining the interface
orientation comes from the connectivity provided by the cell stencil on which
the volume fraction gradient is computed.% \footnote{A good example of an
accurate gradient calculation method on a structured mesh is the Mixed
Youngs-Centered (MYC) gradient calculation method, developed by
\citet{Aulisa2007}. Using a wider stencil to increase the accuracy in
computing the interface normal introduces numerical smoothing of the interface
curvature, as reported by \citet{KotheUnstructured1996}. This is a common issue
of the  geometrical VoF algorithm, however, since the reconstruction error
reduces with the increase of mesh density, the application of local dynamic AMR
significantly reduces both the reconstruction errors and the need for a
reconstruction algorithm of higher order. Additionally, the overall
computational overhead of the solution algorithm is significantly reduced by
reducing the number of cells compared to a static uniformly refined mesh. 

The overall reconstruction algorithm is described by Algorithm
\ref{alg:reconstruction}, defined by two key steps: gradient
calculation and interface positioning. 

\begin{algorithm}[H]
    \centering
    \caption{Reconstruction algorithm}
    \label{alg:reconstruction}

    {\small
      \begin{algorithmic}
       \State compute the gradient of the volume fraction
       \For{cells}
         \If{cell is an interface cell}
             \State initialize the iterative positioning
             \While{interface not positioned}
                 \State secant method: update position
                 \If{secant diverging (two arguments outside search interval)}
                    \While{interface not positioned}
                     \State bisection method: update position 
                    \EndWhile
                 \EndIf
                 \State update secant divergence test (store last argument)
             \EndWhile
         \EndIf
        \EndFor
      \end{algorithmic}
    }
\end{algorithm}


\subsubsection{Gradient calculation}

Calculating the volume fraction gradient to second order of convergence on an
unstructured mesh has been investigated mostly for single phase simulations
(e.g. \citet{lsq-revisiting}) and for volume rendering algorithms in the field
of Computer Graphics (e.g. \citet{IEEE-gradComparison}). When it comes to two-phase 
flows on unstructured mesh, an overview of the possible gradient schemes
is provided by \citet{KotheUnstructured1996}.
\citet{AhnShashkovMultiMaterialRec2007} have used a least squares approach for
gradient calculation on an unstructured mesh for comparing the results with
their MoF reconstruction. 

A steep change of the volume fraction field across the fluid interface within a
single cell layer introduces difficulties in obtaining an accurate gradient
calculation on unstructured meshes. Computation of the gradient on unstructured
meshes based on the approximation stemming from the Gauss divergence theorem is
of second order convergence, however, this gradient scheme does not take into
account point cell neighbors. Therefore, the absolute accuracy of such schemes
deteriorates on fields with large jumps in values, across a narrow band of
cells. In the case of our geometrical VoF algorithm, the reduced absolute
accuracy of the standard Gauss gradient is observed in the form of
instabilities of the reconstructed interface geometry. The instabilities
persist regardless of the choice of the interpolation scheme, since the
interpolation stencil remains unchanged. This is especially true for
tetrahedral and hexahedral unstructured meshes. On polyhedral meshes, each cell
is connected to all neighbor cells by sharing a common face. The face-based
gradient calculation on polyhedral meshes is thus more accurate.

We have used a \emph{Node Averaged Gauss} (NAG) gradient scheme to compute the
orientation of the interface normal. The NAG gradient relies on the
connectivity which is already provided by the unstructured mesh. However, the
increase in absolute accuracy comes from successive interpolation of cell
centered values to mesh points, and averaging of the point values at face
centers. The value of the field is interpolated from cell centers to cell
points using an Inverse Distance Weighting (IDW) interpolation:
%
\begin{equation}
    \phi_p = \sum_{pc} w_{pc} \phi_{pc}
    \label{eq:idwpointphi}
\end{equation}
%
where $pc$ marks cells $c$ surrounding the mesh point $p$ (point-cells), 
$\phi_p$ is the field value which maps to mesh points, and $w_{pc}$ is the
inverse distance weight:
%
\begin{equation}
    w_{pc} = \frac{\frac{1}{||\x_p - \x_{pc}||}}{\sum_{\widetilde{pc}} \frac{1}{||\x_p - \x_{pc}||}},
\end{equation}
%
between the center of the point-cell $pc$ and the mesh point $p$. 

From the mesh point values, the face-centered values are calculated using
\emph{area weighted averaging}. As a first step, the center of the face is
calculated as the algebraic average of the face points, and the initial
face-centered value $\phi_{f}'$ is computed by algebraically averaging the
point values: 
%
\begin{equation} \phi_{f}' = \frac{1}{N_{fp}} \sum_{fp} \phi_{fp},
\end{equation}
%
where $\phi_{f}$ is the \emph{initial face centered value}, $fp$ is a set of
mesh points that comprise a face $f$ (face-points) and ${N_{fp}}$ is the number
of face-points. In a second step, the face is decomposed into triangles using
the face center and points of each face edge ($fp$, $fp+1$). Triangle area is
computed using the vector product its edge vectors, and is then used as a
weighting factor: the final face centered value is then computed as an area
weighted average of the averaged triangle values: 
%
\begin{equation} \phi_f = \sum_t w_t \phi_t, \end{equation}
%
where $w_t$ is the area weight of the triangle:
%
\begin{equation} w_t = \frac{A_t}{\sum_t A_t}, \end{equation}
%
and $\phi_t$ is the average value of the triangle: 
%
\begin{equation} \phi_t = \frac{1}{3} (\phi_f' + \phi_{fp} + \phi_{fp+1}).
\end{equation}
%
As a final step in computing the cell centered gradient, the face centered
field is used to calculate the gradient using the Gauss divergence theorem in
the gradient form: 
%
\begin{equation} 
\grad{\phi}_c \approx \frac{1}{V_c} \sum_f \phi_f \vec{S}_f. 
\end{equation}
%
This algorithm accounts for the volume fraction values stored in \emph{skew
cells} (cells which are connected with the cell in question through a point or
an edge) with their information to be transferred to the interface cell
indirectly, without the explicit computation of the additional connectivity
required by the wide cell stencil on the unstructured mesh. 

The area weighted averaging takes into account the contribution to the final
face-centered value by taking into account the shape of the face. An
alternative calculation of the NAG gradient could involve IDW interpolation
applied to compute the face centered values instead of the area averaging
procedure. This alternative has not been tested since the above averaging 
approach has been used to produce very accurate as well as convergent 
interface reconstruction results.

The \emph{generalized least squares gradient} (generalized LSQ) that supports
arbitrary unstructured meshes as described by \citet{lsq-revisiting} has also
been implemented. For the generalized LSQ gradient, the connectivity stencil
storing the access information to cell point neighbors needs to be
constructed. For a linear least squares gradient, the field is approximated
using the linear part of the Taylor series, i.e. 
%
\begin{equation}
    \phi(\x_{k}) \approx  \phi(\x_{c}) + \grad{\phi}(\x_{c}) (\x_{k} - \x_{c}),
\end{equation}
%
where $c$ and $k$ are the indirect access indices of two cell centers on an
unstructured mesh. Since a wider cell stencil on a three-dimensional mesh
implies a larger number of cell centers than three, a gradient of a scalar
field $\phi_c$ is computed by minimizing the weighted gradient error for each cell $c$:
%
\begin{equation}
    E_c = \sum_k w_{ck}^2 E_{ck}^2,   
    \label{eqn:generalizedLSQminimize}
\end{equation}
%
where the error is defined as (\citep[]{lsq-revisiting})
%
\begin{equation}
    E_{ck} = -[\phi(\x_{c}) - \phi(\x_{k})] + \grad{\phi}(\x_c) (\x_c - \x_k),
\end{equation}
%
and the $w_{ck}$ is the inversed distance weight. The minimization of the error
$E_c$ for a scalar field results with a solution of a  algebraic equation
system for each cell of dimensions $3\times3$, with the dimensions defined by
the three components of the gradient vector. 

\subsubsection{Interface positioning}

An iterative approach to interface positioning
(\citet{RiderReconstructing1998}) using the super-linear convergent Brent's
algorithm (\citet{BrentAlgorithm}) or an analytical approach
(\citet{ScardovelliAnalytical2000}) can be chosen for interface positioning.
\citet{AhnShashkovMultiMaterialRec2007} have shown that the interface
positioning may be accomplished in a simpler way, using a combination of the
secant and bisection iteration methods. The bisection method has guaranteed
linear convergence, and the secant method has super-linear convergence, but may
diverge. 

Figure \ref{fig:positioning} shows the details of our interface positioning
algorithm. A characteristic fill level function is shown which corresponds to
the increase of the volume fraction value as the interface is moved along the
orientation axis given by the gradient of the volume fraction. Because of the
monotonicity of the fill level function we know that the fill level function
has a zero value within the search interval. However, divergence of the secant
method may happen when the iterative procedure reaches the border of the
interval since the fill level function has a \emph{vanishing slope} at the
interval borders ($-\alpha$ or $1 - \alpha$), so the iterative procedure
switches to the bisection method in this case. 

\begin{figure}
    \centering
    %\def\svgwidth{\columnwidth}
    \def\svgwidth{0.5\columnwidth}
    {\footnotesize
        \input{figures/interface-positioning-black-white-scaled.pdf-tex}
    }
    \caption{Geometric Reconstruction: Interface positioning}
    \label{fig:positioning}
\end{figure}

\subsection{Volume fraction advection}

Advection of the volume fraction is based on a geometrical computation of the
\emph{fluxed phase volume} $V_{\alpha}$ on cell faces. This is achieved by
\emph{sweeping} (backtracking) the face along the streamline vectors given by
the interpolated velocity field at the mesh points: 
%
\begin{equation} 
\vec{x}'_{fp} = \vec{x}_{fp} - \vec{U}_{fp}\Delta t,
\label{eq:sweep} 
\end{equation}
%
where $\vec{x}_{fp}$ is the face-point (a mesh point addressed via the cell
face index list).  The swept face, together with the original face, construct
together a so-called \emph{congruent} polyhedron, which may have non-planar faces and
may even be non-convex. The advantage of using unique point velocities lies in
the fact that the \emph{swept polyhedra do not overlap at all}, as shown in
Figure \ref{fig:sweptPolyhedron:sweeping}. However, it has been reported in the
literature \citep{LiovicCVTNA2006} that using the full discrete flow map may
result in highly distorted polyhedra when complex flow fields are present,
introducing errors in mass conservation and numerical boundedness of the volume
fraction field. To the contrary, as we show later in the results part, below the 
expected Courant number limit ($Co<1$) our development of the numerical flow 
solution remains mass conservative and bounded within 
$[0,1]$ using point velocities, even for complex flow fields and both on meshes 
with densities used for standard test cases and on dynamically refined meshes 
using local dynamic AMR. %
We have observed an increased accuracy and robustness of the advection algorithm 
because of the tetrahedral decomposition applied to the computation of the volume 
of the swept polyhedron. %
Analytical expressions for computing the volume of the flux polyhedron rely
exclusively on the convexity of the polyhedron. Using unique point velocities
to sweep the face does introduce some non-planarity of the faces of the swept
as shown for an example shaded non-planar face in Figure
\ref{fig:sweptPolyhedron:sweeping}. However, the distortion of the polyhedron
will always be small enough, such that the tetrahedral decomposition removes the
face non-planarity and delivers accurate results for the volume of the swept
polyhedron $S_f$ with non-planar faces, which allows for an efficient iterative
flux correction procedure that completely avoids topological operations on the
swept polyhedron. 

Figure \ref{fig:sweptPolyhedron:sweeping} shows the construction of the swept
polyhedron: face point velocities $\U_{fp}$ are interpolated using 
\eqref{eq:idwpointphi} and multiplied by $-\Delta t$ to track the
face back (Equation \ref{eq:sweep}). Tetrahedral decomposition is then used to
compute the volume of the volumetric flux per each face $f$ and the polyhedron
face is swept along the stream vectors iteratively (Figure
\ref{fig:sweptPolyhedron:correcting}) until the volume of the swept polyhedron
becomes equal to the volume given by the volumetric flux $F_f$: 
%
\begin{equation}
    V_{f,F} = \| F_f \| \Delta t,
    \label{eq:VF}
\end{equation}
converging into the final \emph{flux polyhedron} used to compute the volume of
the fluxed phase accross face $f$.
%
\begin{figure}
   \begin{subfigure}[b]{0.45\columnwidth}
       \centering
       \def\svgwidth{\columnwidth}
       {\footnotesize
            \input{figures/swept-polyhedron-black-white-scaled.pdf-tex}
       }
       \caption{Sweeping faces with point velocities produces non-overlapping swept polyhedra} 
       \label{fig:sweptPolyhedron:sweeping}
   \end{subfigure}
   \begin{subfigure}[b]{0.45\columnwidth}
       \centering
       \def\svgwidth{\columnwidth}
       {\footnotesize
            \input{figures/swept-polyhedron-black-white-scaled-iterative.pdf-tex}
       }
       \caption{Iterative correction of the swept polyhedron to match the volume given by the volumetric flux}
       \label{fig:sweptPolyhedron:correcting}
   \end{subfigure}
\end{figure}
%
The velocity field obtained at the face points using inversed distance
interpolation is not divergence free in the discrete sense, thus a divergence 
correction becomes necessary. For this, we define the volume of the face-swept 
flux polyhedron:
%
\begin{equation}
    V_{f,S} = \sum_{f,T} V_{f,T}, 
\end{equation}
%
where $V_{f,T}$ is the volume of a tetrahedron produced as a result of
tetrahedral decomposition of the swept-face polyhedron $S_f$ generated for each
face $f$ of the mesh (Figure \ref{fig:sweptPolyhedron:sweeping}).  Using the volume
of the decomposed face-swept flux polyhedron, the discrete swept volumetric
flux can be computed as
\begin{equation}
    F_{f,S} = \frac{V_{f,S}}{\Delta t}.
\end{equation}
As already noted, this volumetric flux is not expexted to satisfy the discrete 
divergence free condition \eqref{eq:discretedivergence} necessary for the mass 
(volume) conservation of a pseudo-staggered solution algorithm, i.e.\
\begin{equation}
    \sum_f {F_{f,S}} \sgn(F_f) \ne 0,
\end{equation}
where $\sgn(F_f)$ is the sign function of the volumetric flux: 
\begin{equation} 
   \sgn(F_f) = \begin{cases} 
                    +1 , F_f > 0 \\
                    -1 , F_f < 0  \cdot  
              \end{cases},
\end{equation}
%
%Since the velocity field in the face points is not divergence free in a discrete sense
%
which needs to be computed once per face since the volume of the swept-face polyhedron 
will always be positive. Subsequently a \emph{divergence correction} of the swept 
polyhedrons is to be performed to ensure mass conservation. %
This practice has proven sufficient for the computation of the fluxed phase volumes, 
since the point velocity field is then only needed to provide the direction for face 
sweeping (streamline direction).%

In literature, there are three approaches commonly used for the divergence
correction: a scalar coefficient correction \citep{LiovicCVTNA2006}, analytical
correction \citep{HernandezLopezPart1Advection2008} and a parametric correction
\citep{MencingerZunAdaptive2011}. We have developed a novel efficient and 
iterative approach: %
the \emph{swept polyhedron} computed by sweeping the face is corrected iteratively 
until its volume is equal to the volume given by the fluxed volume defined by the 
scalar volumetric flux (Equation \ref{eq:VF}).
The super-linear convergent secant method is applied to \emph{move the end points
of the flux polyhedron} until its volume corresponds to the volumetric flux volume. 
This algorithm does not involve topological changes of the flux polyhedron and the 
method converges in a very small number of steps (3-5 on average, evaluated for 
the shear flow test case) to machine tolerance since the initial volume of the flux 
polyhedron is in magnitude near to the magnitude of the volumetric flux, and the 
target function is a fill level function with a monotonic increase. Avoiding 
topological operations in the correction of the flux polyhedron increases the speed 
of the correction significantly. We have tested the algorithm for different shapes 
of swept faces, applying random perturbations in the orientation of swept edges 
and have found the correction algorithm to be robust and accurate. Even with 
random perturbations of swept point coordinates in magnitudes up to 20\% of the 
characteristic length of the face (Figure \ref{fig:sweptPolyhedron:sweeping}), 
the calculation of the volume magnitude remains correct to machine tolerance, 
allowing for the iterative correction algorithm to converge very quickly. 

%The reason behind such stability of the algorithm when dealing with non-planar
%faces lies in the way the polyhedral cells are handled: cells are considered
%to be \emph{logical} polyhedrons.  Logical polyhedrons do not impose planarity
%of the faces, and the unstructured mesh built from logical polyhedrons does
%not directly require the cells to be convex. The computation of the flux
%polyhedron volume is thus based on the tetrahedral decomposition, which works
%very well even for moderately non convex polyhedrons with non-planar faces.

\begin{figure}
   %\begin{subfigure}[b]{0.5\columnwidth}
       \centering
       \def\svgwidth{\columnwidth}
       {\footnotesize
           \input{figures/polyhedral-flux-intersection-black-white-scaled.pdf-tex}
       }
       %\caption{Intersecting the corrected flux (swept) polyhedron and PLIC plane} 
       %\label{fig:volPhaseFlux:intersection}
   %\end{subfigure}
   %\begin{subfigure}[b]{0.5\columnwidth}
       %\centering
       %\def\svgwidth{\columnwidth}
       %{\footnotesize
           %\input{figures/polyhedral-flux-final-black-white-scaled.pdf-tex}
       %}
       %\caption{Result of intersection: fluxed phase volume contribution}
       %\label{fig:volPhaseFlux:contribution}
   %\end{subfigure}

    \caption{Computing the volume of the fluxed phase using intersections.} 
    \label{fig:volPhaseFlux}
\end{figure}

After the swept polyhedrons have been corrected, the contributions of the volumetric
phase flux coming from the surrounding cells needs to be computed. Obtaining a
flux contribution using geometrical intersections is shown schematically in
Figure \ref{fig:volPhaseFlux} for a single face. The corrected swept 
polyhedron $S_f$ (dotted lines) is initially clipped with each cell of the 
advection stencil:
%
\begin{equation}
   S'_f = S_f \cap_s C_s,  
\end{equation}
%
where $C_s$ is the cell of the chosen flux stencil $s$. The flux stencil $s$
needs to incorporate point and edge face neighbors, otherwise those
contributions to the total fluxed phase volume will be missing and interface
motion will be artificially distorted in those directions, as is the case
for flux-based methods naively utilizing algebraic advection schemes on 
unstructured hexahedral meshes. In case the candidate cell is an interface 
cell, a subsequent intersection is performed with the positive halfspace 
$I^+$ defined by the PLIC plane (grey colored polygon):
%
\begin{equation}
    S''_f = S'_f \cap I^+.
\end{equation}
%
This is the final face contribution to the total fluxed phase volume as illustrated
in Figure \ref{fig:volPhaseFlux}. The total value of the volumetric flux is
then computed as the sum of the fluxed phase volume contributions calculated for
all cells of the face flux stencil. The calculation of the total fluxed phase
volume across a face is performed for an \emph{outflow cell} only using 
\eqref{equation:vol-phase-flux}, resulting in a final decrease of the volume 
fraction for this cell, and the corresponding increase for the \emph{inflow cell}. 

We use a \emph{narrow band} of interface cells for the advection of the volume
fraction field. %Using non divergence-free mesh point velocity field throughout
%the domain would introduce errors in mass conservation. Moreover, 
Since geometrical intersections are needed for the geometrical reconstruction and
advection steps, localizing the computation to a narrow band of cells strongly
increases the efficiency of the overall algorithm. In case of a complex
velocity field and a naive advection operating on the whole domain,
\emph{wisps} appear in the bulk of both phases. Hence, locating and removing
wisps is greatly simplified for a \emph{narrow band} advection since their
initial location is always in the bulk cells that are direct neighbors to
interface cells. To remove the wisps, a redistribution algorithm of
\citet{HarvieFletcherStream2000} has been implemented, with applied
modifications that are necessary for fast computation of large stencils on
arbitrary unstructured meshes. Both void and phase wisps are redistributed
conservatively to the surrounding interface cells, under the constraint that
the removal of mass or its addition is not allowed to cause further creation of
over/undershoots or wisps. 

The narrow cell band is constructed using neighborhood cell information and the
volume fraction field, i.e. the algorithmic task for the advection is: advect
only in interface cells and their first \emph{point neighbors}. To avoid
computing additional point-cell neighbor connectivity which would then be
updated during the evolution of the interface, we have re-used the gradient of
the volume fraction field as an indicator for the bulk cells that are not
involved in the advection. This is schematically shown in Figure
\ref{fig:parallel:gradZeroCondition}.  A bulk cell is  characterized as a cell
for which the gradient of the volume fraction is zero.

Since only the outflow from cells is calculated, the change of the volume
fraction is not applied for cells that lie in the bulk of the phases.  The
fluxed phase volume is computed also for cell faces when the outflow is
present for the bulk cell. However, in this case the update of the volume
fraction field is not executed for the bulk cell. The re-use of the gradient
field as a criterion for constructing the narrow band of cells involved in the
advection significantly increases the efficiency of the method, since the
gradient is already computed as a requirement of the reconstruction algorithm.

\subsection{Parallel implementation}

The parallel implementation of the reconstruction algorithm is simplified when
the NAG gradient is utilized for the computation of the interface normal, since
no additional communication structures need to be prepared and communicated
over the process boundary. In the case of the generalized LSQ algorithm,
computing the additional connectivity involves multiple issues: parallel
communication involving more complex stencil structures, update of the stencil
structures when topological changes are applied to the mesh (AMR), and caching
of the additional connectivity to increase algorithm efficiency. Since the NAG
gradient has provided sufficiently accurate results, we have left the parallel
implementation of the generalized LSQ gradient as future work if proved to be
necessary.

\begin{figure}
   \begin{subfigure}[b]{0.45\columnwidth}
       \centering
       \def\svgwidth{0.6\columnwidth}
       {\footnotesize
           \input{figures/advectionGradientSimpleFront.pdf-tex}
       }
       \caption{Computing a narrow band of cells} 
       \label{fig:parallel:gradZeroCondition}
   \end{subfigure}
   \begin{subfigure}[b]{0.45\columnwidth}
       \centering
       \def\svgwidth{0.6\columnwidth}
       {\footnotesize
           \input{figures/advectionGradientParallelComm.pdf-tex}
       }   
       \caption{Communicating the outflow flux fields}
       \label{fig:parallel:communication}
   \end{subfigure}

    \caption{Parallel implementation: computing the fluxed phase volume}    
    \label{fig:parallel:sphereFine}
\end{figure}

Using the gradient to determine the bulk cells simplifies the parallel
implementation of the advection algorithm as shown schematically on Figure
\ref{fig:parallel:gradZeroCondition}. Additional connectivity required to
identify the cell candidates would otherwise be communicated across process
boundaries, as well as the geometrical information of the mesh cells. However,
when the volume fraction gradient is used to determine the narrow band of
cells, there is no need for additional connectivity. Since the outflow field is
computed for each separate domain, the boundary outflow fluxes need to be
exchanged (simultaneously) across process boundaries. The outflow from one
domain represents the inflow field for the other domain, as shown in Figure
\ref{fig:parallel:communication}. In order to avoid \emph{deadlocking}
resulting from simultaneous communication, the parallel code utilizes \emph{non
blocking} communications \citep{OpenMPI}. 
